\indent ಪ್ರಸ್ತಾರ ಹಾಕುವುದು : ಯಾವುದೇ ಪದ್ಯದ ಪ್ರತಿಯೊಂದು ಸಾಲಿನಲ್ಲಿರುವ ಅಕ್ಷರಗಳಿಗೆ ಗುರು-ಲಘುಗಳನ್ನು ಗುರುತು ಮಾಡುವ ಕ್ರಮವನ್ನು ಪ್ರಸ್ತಾರ ಹಾಕುವುದು ಎನ್ನುತ್ತೇವೆ. ಪ್ರಸ್ತಾರ ಹಾಕಿದ ನಂತರ ಗುರುವಿಗೆ ಎರಡು ಮಾತ್ರೆ ಮತ್ತು ಲಘುವಿಗೆ ಒಂದು ಮಾತ್ರೆ-ಈ ಪ್ರಕಾರವಾಗಿ ಎಣಿಕೆ ಮಾಡಿ ಮೂರು ಮಾತ್ರೆ, ನಾಲ್ಕು ಮಾತ್ರೆ ಅಥವಾ ಐದು ಮಾತ್ರೆಯ ಗುಂಪುಗಳನ್ನು ವಿಂಗಡಿಸಬೇಕು; ಈ ಗುಂಪುಗಳಿಗೆ “ಮಾತ್ರಾಗಣ” ಗಳೆಂದು ಹೆಸರು.\par

\kbvsechead{ಮಾತ್ರಾಗಣದ ಪದ್ಯಗಳ ಲಕ್ಷಣ :}

\vspace{4pt plus 4pt minus 1pt}\noindent{\kbvsecfont\fontsize{11pt}{14.5pt}\selectfont\color{kbvaccentdark} (೧)\hspace{0.45em}ಕಂದ:}\hspace{0.5em}ಕಂದ ಪದ್ಯದಲ್ಲಿ ನಾಲ್ಕು ಪಾದ (ಸಾಲು) ಗಳಿರುವುವು. ಇವುಗಳ ಪೈಕಿ ಒಂದು ಮತ್ತು ಮೂರನೆಯ ಪಾದಗಳಲ್ಲಿ ನಾಲ್ಕು ಮಾತ್ರೆಯ ಮೂರು ಗಣಗಳಿರುವುವು; ಎರಡು ಮತ್ತು ನಾಲ್ಕನೆಯ ಪಾದಗಳಲ್ಲಿ ನಾಲ್ಕು ಮಾತ್ರೆಯ ಐದು ಗಣಗಳಿರುವುವು.\par

\vspace{3pt plus 2pt}\noindent\hspace*{1.2em}\kbvlabel{ಉದಾ}\par\vspace{1pt}

\indent \laghu{}\laghu{}\laghu{} \laghu{} \laghu{}\laghu{}\laghu{} \laghu{} \laghu{}\laghu{}\guru{}\par

\indent 4 \ensuremath{\times} 3....ಭವದನು ಜನರುಣ ಜಲಮಂ \laghu{}\laghu{}\guru{} \laghu{}\laghu{}\guru{} \laghu{}\guru{}\laghu{} \laghu{}\laghu{} \laghu{}\laghu{}\guru{}\par

\begin{kbvverse}
\noindent 4 \ensuremath{\times} 5....ಸವಿನೋ ಡಿದನಿಂ ತು ನಿನ್ನ ಬಲಜಲ ನಿಧಿಯಂ ।
\end{kbvverse}

\indent \laghu{}\laghu{}\guru{} \laghu{}\laghu{}\guru{} \laghu{} \laghu{}\guru{}\par

\indent 4 \ensuremath{\times} 3....ಸವಿನೋ ಡಿದೆನೀ ಕೊ ಳನಂ \laghu{}\laghu{}\guru{} \laghu{}\guru{}\laghu{}\laghu{} ೧ \guru{}\laghu{} \laghu{}\laghu{}\guru{} \guru{} \guru{}\par

\begin{kbvverse}
\noindent 4 \ensuremath{\times} 5....ತವೆಪೀ ರ್ದುಬಳಿ ಕೈ ನಿನ್ನ ಸನಿಯಂ ನೋಳ್ಪಂ ॥
\end{kbvverse}

\indent ವಿಶ್ಲೇಷಣೆ : (1) ಈ ಪದ್ಯದಲ್ಲಿ ನಾಲ್ಕು ಪಾದಗಳಿವೆ. (2) ಒಂದು ಮತ್ತು ಮೂರನೆಯ ಪಾದಗಳಲ್ಲಿ ನಾಲ್ಕು ಮಾತ್ರೆಯ ಮೂರು ಗಣಗಳಿವೆ. (3) ಎರಡು ಮತ್ತು ನಾಲ್ಕನೆಯ ಪಾದಗಳಲ್ಲಿ ನಾಲ್ಕು ಮಾತ್ರೆಯ ಐದು ಗಣಗಳಿವೆ. (4) ಪೂರ್ವಾರ್ಧ ಮತ್ತು ಉತ್ತರಾರ್ಧಗಳಲ್ಲಿ ಕ್ರಮವಾಗಿ ತಲಾ ಎಂಟು ಮಾತ್ರಾಗಣಗಳಿವೆ. (5) ಪೂರ್ವಾರ್ಧ ಹಾಗೂ ಉತ್ತರಾರ್ಧದ ಮಾತ್ರಾಗಣಗಳ ವಿಷಮ ಸ್ಥಾನಗಳಲ್ಲಿ ಮಧ್ಯಗುರುವುಳ್ಳ “\laghu{}\guru{}\laghu{}” ಈ.\par
